\section{Trabajo Adelantado}\label{sec:trabajo_adelantado}

El trabajo adelantado de esta memoria se ha orientado a reducir el riesgo técnico del proyecto antes de intervenir directamente el algoritmo asociado a \texttt{ApplicativeDo} en \textsf{GHC}. En lugar de limitarse a una revisión bibliográfica o a una caracterización puramente conceptual del problema, este avance se concretó en un repositorio experimental funcional \cite{RomeroApplicativeDoRepo}, reproducible y documentado, preparado específicamente para estudiar el pipeline real del compilador y habilitar modificaciones controladas sobre su implementación. En consecuencia, el trabajo adelantado no solo acredita familiaridad con el problema de investigación, sino que demuestra que ya existe una base técnica concreta sobre la cual desarrollar la propuesta de memoria.

\subsection{Infraestructura reproducible y baseline experimental}

Un primer avance fundamental fue la consolidación de un entorno reproducible basado en un contenedor \textsf{Docker} de desarrollo (\textit{Dev-Container}), con el objetivo de eliminar variaciones del entorno local y asegurar que el proyecto pudiera reconstruirse de manera consistente. Esta decisión fue necesaria porque el trabajo sobre \textsf{GHC} depende de una cadena de herramientas precisas, y pequeñas diferencias de versiones pueden impedir la compilación del compilador o introducir errores difíciles de diagnosticar. En este contexto, se resolvieron incompatibilidades de versiones asociadas a \textit{ghcup}, particularmente problemas de disponibilidad de versiones, y se fijó una combinación funcional de herramientas para el entorno experimental: \textsf{GHC} de bootstrap \texttt{9.14.1}, \texttt{cabal-install} \texttt{3.14.2.0}, \texttt{alex} \texttt{3.5.4.0} y \texttt{happy} \texttt{2.2}. Adicionalmente, se incorporaron las dependencias del sistema necesarias para ejecutar \textit{./boot} en \textsf{GHC}, en particular \texttt{autoconf}, \texttt{automake} y \texttt{libtool}, y se añadieron scripts de verificación de la toolchain dentro del directorio \textit{.devcontainer}, de modo que la validación del entorno quede automatizada y no dependa de inspección manual.

En paralelo, se estabilizó la base del compilador utilizada para la experimentación. Inicialmente, se partía desde un commit antiguo de \textsf{GHC} con problemas de reproducibilidad debido a submódulos históricos rotos u obsoletos, lo que hacía riesgosa cualquier intervención posterior. Para superar este problema, se migró a un baseline moderno del submódulo \textsf{GHC}, fijando como commit de referencia \textit{0b36e96cb93db71f201aaa055c4a90b75a8110ef} por su mayor actualidad y no el commit específico asociado a la extensión \texttt{ApplicativeDo} (2015) \cite{ghc_commit_8ecf6d8}. Sobre esta base se crearon y ajustaron scripts para el alta, inicialización y verificación del submódulo \textit{vendor/ghc}, así como de sus submódulos anidados, y se incorporó automatización mediante un archivo \textit{Makefile} para simplificar el proceso de setup y reproducción. Más importante aún, se tomó la decisión de privilegiar una estrategia de patches sobre el submódulo en lugar de modificar un fork sin trazabilidad, lo que fortalece simultáneamente la reproducibilidad, la mantenibilidad y la trazabilidad de los cambios locales.

\subsection{Construcción de GHC y uso controlado del compilador del experimento}

Un segundo frente de trabajo consistió en documentar y validar el flujo real de compilación de \textsf{GHC} dentro de \textit{vendor/ghc}. En particular, se dejó establecido el pipeline efectivo de construcción del compilador mediante las etapas \textit{./boot}, \textit{./configure} y \textit{./hadrian/build -j"\$(nproc)" --flavour=devel2}. Esto fue necesario para pasar desde una caracterización teórica del código fuente a una base experimental verificable, donde el compilador pudiera efectivamente reconstruirse y utilizarse como herramienta de trabajo. Junto con ello, se verificó el uso de los binarios generados en las carpetas  \textit{stage1} y \textit{stage2} dentro del submódulo, lo que permitió confirmar que el flujo de build no solo finaliza correctamente, sino que produce artefactos ejecutables aptos para la experimentación.

Sobre esta base, se habilitó explícitamente la compilación de programas del experimento usando \textsf{GHC} construido desde el submódulo, en vez del \textsf{GHC} global del sistema. Esta decisión es crucial para la validez experimental de la memoria: si las pruebas se ejecutaran con el compilador global, cualquier modificación realizada dentro de \textit{vendor/ghc} dejaría de ser observable en el experimento, comprometiendo la interpretación de resultados. Al usar de forma controlada el compilador reconstruido desde el submódulo, el repositorio experimental garantiza que la observación del comportamiento de \texttt{ApplicativeDo} corresponde realmente a la implementación intervenida. En otras palabras, este avance no solo habilita la compilación de \textsf{GHC}, sino que asegura que las etapas posteriores de la memoria trabajen sobre una base experimental fiel al código efectivamente modificado.

\subsection{Plataforma experimental para observar \texttt{}{ApplicativeDo}}

Sobre la infraestructura anterior se desarrolló y refinó una plataforma experimental específica en \textit{experiments/applicative-do/}, destinada a observar de forma directa el comportamiento del compilador frente a distintos patrones de uso de \texttt{ApplicativeDo}. El archivo \textit{Main.hs} del experimento fue ampliado hacia un caso más representativo y compacto, incorporando escenarios que ejercitan segmentos applicative puros, dependencias monádicas, BodyStmt, patrones estrictos, patrones refutables, bloques anidados y casos que fuerzan el uso de \texttt{join}. Esta ampliación fue necesaria para que el experimento no se limitara a ejemplos mínimos, sino que cubriera una variedad de situaciones estructuralmente relevantes para el estudio del algoritmo.

Asimismo, se implementó \textit{run.sh} para compilar y ejecutar el experimento de forma reproducible, con autodetección del compilador \textsf{GHC} construido desde el submódulo y soporte para la variable \textit{GHC\_BIN}. A esto se sumó la incorporación de flags orientados a la observabilidad del pipeline del compilador: \textit{DUMP\_DS=1} para obtener el dump de desugaring, \textit{DUMP\_PIPELINE=1} para generar dumps del renamer, typechecker y desugar (\textit{-ddump-rn}, \textit{-ddump-tc}, \textit{-ddump-ds}), \textit{DUMP\_SIMPL=1} para obtener el Core simplificado, y \textit{EXTRA\_GHC\_FLAGS="..."} para extender la observación con opciones adicionales. También se agregó un manifiesto de dumps generados en \textit{build/dump-manifest.txt}. El valor de esta plataforma es doble: por una parte, permite explorar de manera controlada el proceso de \textbf{Rearrangement} y \textbf{Desugar} asociado a \texttt{ApplicativeDo}; por otra, constituye una base práctica para contrastar futuras modificaciones del algoritmo contra ejemplos concretos y artefactos verificables del compilador.

\newpage

\subsection{Documentación técnica, trazabilidad e impacto en la viabilidad}

Un último bloque de trabajo adelantado fue la consolidación de documentación técnica específica, orientada no solo a describir el repositorio, sino a capturar conocimiento operativo y arquitectónico que de otro modo quedaría disperso o implícito. En particular, se actualizaron \textit{README.md}, \textit{docs/toolchain.md}, \textit{docs/architecture.md} y \textit{experiments/applicative-do/README.md}. Dentro de este conjunto, \textit{docs/architecture.md} tiene un valor especialmente alto para la memoria, pues deja trazado un mapa exhaustivo de \texttt{ApplicativeDo} en \textsf{GHC} actual: se describe el recorrido Parser -> \ Renamer -> \ Typechecker -> \ Zonk -> \ Desugar, las estructuras \texttt{ApplicativeStmt} y \texttt{ApplicativeArg}, funciones críticas como \texttt{rearrangeForApplicativeDo}, \texttt{segments}, \texttt{mkApplicativeStmt}, \texttt{tcApplicativeStmts}, \texttt{dsDo}, entre otras, así como los puntos de extensión más relevantes para experimentación. Esto transforma un conocimiento que normalmente requeriría una exploración larga y costosa del código fuente en un activo reutilizable y verificable.

En conjunto, estos avances permiten afirmar que el proyecto ya dispone de una base experimental real sobre la cual construir la investigación propuesta. Se cuenta con un entorno reproducible y validado, con una estrategia clara de mantenimiento sobre \textsf{GHC}, con un baseline moderno y verificable, con una ruta documentada de compilación del compilador, con una plataforma experimental que permite observar \texttt{ApplicativeDo} en funcionamiento y con documentación técnica suficiente para localizar los puntos donde la memoria deberá intervenir. Desde el punto de vista metodológico, esto mejora la reproducibilidad, fortalece la trazabilidad, incrementa la capacidad experimental y reduce de manera importante el riesgo técnico del proyecto. En consecuencia, la memoria puede concentrarse en su núcleo de investigación, es decir, la generación de permutaciones válidas de statements, su integración al algoritmo base y el estudio de su efecto sobre los planes de ejecución, todo esto sobre una base robusta y verificable.




\newpage
